\begin{table}[htbp]
  \centering
  \caption{关键Tamarin验证结果}
  \label{tab:key-verification-results}
  \small
  \begin{tabularx}{\linewidth}{@{}L{0.08\linewidth}L{0.40\linewidth}L{0.10\linewidth}L{0.12\linewidth}Y@{}}
    \toprule
    配置 & 验证性质 & 类型 & 结果 & 论证作用 \\
    \midrule
    放宽 & 一个精确来源支持两个同批接受\newline
      {\footnotesize\path{one_send_two_accepts_exists}} & 存在迹 & 已验证 & 重复接受见证 \\
    放宽 & 作用域内发生注入性\newline
      {\footnotesize\path{receiver_accept_injective}} & 全迹 & 被反例否定 & 对应全迹性质失败 \\
    \midrule
    消息级 & 消息区分\newline
      {\footnotesize\path{accepted_batch_has_distinct_messages}} & 全迹 & 已验证 & 维护消息坐标性质 \\
    消息级 & 作用域内发生注入性\newline
      {\footnotesize\path{receiver_accept_injective}} & 全迹 & 已验证 & 排除精确重复形状 \\
    消息级 & 同一参与方不同消息批次可达\newline
      {\footnotesize\path{same_party_different_messages_batch_exists}} & 存在迹 & 已验证 & 支撑非替代性结论 \\
    \midrule
    参与方级 & 参与方区分\newline
      {\footnotesize\path{accepted_batch_has_distinct_parties}} & 全迹 & 已验证 & 维持目标参与方关系 \\
    参与方级 & 有效不同参与方批次可达\newline
      {\footnotesize\path{distinct_party_batch_exists}} & 存在迹 & 已验证 & 排除真空成立 \\
    \bottomrule
  \end{tabularx}
  \par\smallskip
  \begin{minipage}{\linewidth}
    \footnotesize 注：存在迹对应exists-trace，全迹对应all-traces；“已验证”对应verified，“被反例否定”对应falsified。性质名下方的小号标识用于追踪对应lemma。
  \end{minipage}
\end{table}
